% -----------------------------------------------------------------------------
% Section 6: Conclusion
% Paper 5 — Constitutional Code (Separation of Powers)
% Security audit patch applied 2026-03-29:
%   Patch 5.1: Open problem paragraph on formal composition (before Lessig)
% -----------------------------------------------------------------------------

\section{Conclusion}
\label{sec:conclusion}

We have shown that the separation of powers---legislative, executive,
judicial---can be encoded directly in a type system, producing
guarantees that are structural rather than conventional. The
Constitutional Completeness Corollary establishes that the three
branches are capability-disjoint: no program that satisfies the BTV
type laws can simultaneously legislate policy, execute decisions, and
adjudicate disputes. This is not a policy recommendation; it is a
theorem.

The result builds on Papers~1 and~2 by adding a third invariant: not
only must every decision be evidenced (Paper~1) and persisted
(Paper~2), but the actors who set policy, execute decisions, and audit
outcomes must be structurally prevented from usurping each other's
roles. Together, the three papers provide a compositional account of
algorithmic accountability at compile-time, delivery-time, and
separation-of-powers level.

\paragraph{Open problem: formal composition.}
The Constitutional Completeness corollary
(Corollary~\ref{cor:completeness}) proves that the three branches
are capability-disjoint in a single BTV institution. What it does
not prove is that the individual security theorems of Papers~1--3
\emph{compose} into a single end-to-end guarantee without
introducing new attack surfaces at the layer boundaries. Paper~3
(Section~7.3 of~\cite{btv-ar2026}) states a composition conjecture;
proving it---ideally in a proof assistant such as Lean~4---is the
most important formal task remaining in the series. Until this
theorem is established, the end-to-end guarantee rests on an
informal layering argument: the three predicates operate at
compile-time, delivery-time, and audit-time respectively, with
non-overlapping trust assumptions. We believe this argument is
sound, but acknowledge that it has not been machine-checked.

Lessig wrote that ``code is law''; we argue the converse is equally
true and more actionable: law is code, and its invariants can be
enforced by a compiler. The accountability crisis in AI governance is
not primarily a problem of bad intentions but of missing enforcement
mechanisms. Linear types, transparency logs, and ZK proofs provide
those mechanisms---not as optional best practices but as compile-time
constraints that make accountability violations impossible to express.
